\chapter{Autism Treatment}

\section*{Definition and Overview}
Treatment for autism is not about curing autism or making an autistic person "normal"; autism is a lifelong pattern of development, and the goal of support is to build skills, reduce distress and improve quality of life. "Treatment" therefore means a set of therapies, educational supports and environmental adjustments tailored to the person's needs and strengths. The most effective approaches begin early, involve the family, respect the person's autonomy and are delivered by trained professionals. Alongside developmental therapies, medicines may be used for specific coexisting problems, and families need reliable guidance to avoid ineffective or harmful "cures" that are sometimes promoted.

\section*{Epidemiology}
Because treatment is individualised, the pattern of services used varies widely. Most children diagnosed with ASD in Australia access early intervention therapy, and many are funded through the National Disability Insurance Scheme (NDIS). Speech pathology and occupational therapy are the most commonly used therapies, often alongside behavioural and psychological support. Coexisting conditions that may need treatment are common: roughly half of autistic people have ADHD symptoms, and anxiety and depression affect many adolescents and adults. Access to services varies by location and income, and waiting lists for assessment and therapy are common, particularly outside major cities.

\section*{Aetiology and Pathophysiology}
Treatment is based on how autism affects learning, communication and daily function. The brain differences seen in autism mean that skills such as understanding spoken language, reading social cues and adapting to change are learned differently, so support must be concrete, structured and repeated.

\begin{itemize}
    \item \textbf{Learning style:} many autistic people learn best with visual information and predictable routines
    \item \textbf{Communication:} therapies may target spoken language or introduce alternative ways to communicate
    \item \textbf{Sensory processing:} differences in how the senses are experienced drive the need for sensory-friendly environments
    \item \textbf{Coexisting conditions:} anxiety, ADHD, epilepsy and sleep problems each need their own treatment
    \item \textbf{No medication for autism itself:} medicines target specific symptoms or coexisting disorders
\end{itemize}

\section*{Clinical Features}
The features that treatment seeks to support fall into several practical areas of daily life.

\begin{itemize}
    \item Communication: from non-speaking or delayed speech to difficulties with conversation and social timing
    \item Everyday skills: dressing, eating, toileting, sleep and household routines
    \item Sensory difficulties: overload from noise, light or touch, or seeking out particular sensations
    \item Repetitive behaviours and insistence on routine, which may bring comfort or become distressing
    \item Anxiety, meltdowns or behavioural distress when needs are unmet
    \item Social participation: friendships, school, work and community activities
    \item Individual strengths: memory, focus, honesty and special interests that can be developed
\end{itemize}

\section*{Differential Diagnosis}
Choosing the right treatment starts with making sure the diagnosis is correct and that coexisting conditions are not missed.

\begin{itemize}
    \item \textbf{Intellectual disability:} needs educational and functional support, with or without autism
    \item \textbf{ADHD:} often coexists; stimulant or other medicines may help if symptoms are severe
    \item \textbf{Anxiety and depression:} common and treatable, but easily overlooked in autistic people
    \item \textbf{Epilepsy:} needs specialist review and seizure medicines in the affected minority
    \item \textbf{Sleep disorders:} behavioural programmes and sometimes medicines improve sleep
    \item \textbf{Language disorders:} need speech therapy even when autism is not the cause
    \item \textbf{Trauma or abuse:} safeguarding concerns require a different response
\end{itemize}

\section*{When to Seek Medical Care}
See a GP or paediatrician if a child is not meeting developmental milestones, if behaviour is very difficult to manage, if there are signs of epilepsy or serious sleep problems, or if there is concern about vision or hearing. Seek a review if the current therapy plan is not helping, or if new symptoms such as low mood, self-injury or weight change appear.

\textbf{Call 000 (or local emergency services) for:}
\begin{itemize}
    \item A seizure lasting more than a couple of minutes
    \item Serious self-injury or an injury needing urgent care
    \item A child or adult who has wandered and is missing
    \item A mental health crisis with immediate risk to safety
\end{itemize}

\section*{Diagnosis and Investigations}
Before treatment is planned, a thorough assessment establishes the diagnosis and the person's specific support needs.

\begin{itemize}
    \item \textbf{Diagnostic assessment:} by a paediatrician, psychiatrist or clinical psychologist using standardised tools
    \item \textbf{Speech and language assessment:} to plan communication therapy
    \item \textbf{Occupational therapy assessment:} for sensory and daily-living skills
    \item \textbf{Psychological assessment:} of learning, memory and emotion
    \item \textbf{Medical review:} for epilepsy, sleep, hearing and vision
    \item \textbf{Review of coexisting conditions:} screening for anxiety, depression and ADHD
    \item \textbf{NDIS planning:} linking the assessment to funded supports
\end{itemize}

\section*{Management}
Good autism support is a team effort, coordinated across therapy, family, school and health services.

\begin{itemize}
    \item \textbf{Early intervention:} structured programmes that build communication, play and daily skills in young children
    \item \textbf{Behavioural therapy:} positive, evidence-based strategies that teach skills and reduce distress when used with trained therapists
    \item \textbf{Speech pathology:} developing spoken language and alternative and augmentative communication
    \item \textbf{Occupational therapy:} sensory regulation, motor skills and everyday independence
    \item \textbf{Psychology:} help with anxiety, emotion regulation and social understanding
    \item \textbf{Education supports:} school adjustments, individual plans and aide support
    \item \textbf{Family coaching and respite:} supporting parents to understand and respond to their child
    \item \textbf{Medicines:} for specific coexisting problems, such as severe anxiety, depression, ADHD or epilepsy, or for marked irritability when it causes distress
    \item \textbf{Transition support:} preparing for secondary school, work and independent living
    \item \textbf{Adult supports:} employment, housing and community programmes
\end{itemize}

\section*{Complications}
The main risks come not from autism itself but from distress that is not recognised, and from treatments that are unproven.

\begin{itemize}
    \item Anxiety, depression and burnout, especially in adolescence and adulthood
    \item Behavioural crisis or self-injury when communication and sensory needs are unmet
    \item Missed or delayed treatment of coexisting conditions such as epilepsy or ADHD
    \item School exclusion or difficulty finding work without adequate support
    \item Family exhaustion and financial strain from the cost of therapies
    \item Harm from unproven "treatments", including restrictive diets, chelation or other unsafe practices
    \item Over-reliance on therapies with little evidence, delaying access to effective support
\end{itemize}

\section*{Prognosis and Outcomes}
Outcomes depend on many factors, but early, consistent, well-matched support is associated with better communication, greater independence and improved wellbeing. Many children who receive good early intervention develop spoken language and attend mainstream school, and many autistic adults live independently and work. People with intellectual disability or more complex needs may continue to require substantial support. Medicines help some people significantly for coexisting conditions, while not changing the underlying autism. Overall, acceptance, predictability and good support tend to predict the best quality of life.

\section*{Prevention and Self-Care}
There is no way to prevent autism, and no evidence that any diet, medicine or procedure can cure it. The priorities are to avoid harm, support wellbeing and make everyday life work well.

\begin{itemize}
    \item Seek early assessment and evidence-based support
    \item Build calm, predictable routines at home and school
    \item Provide sensory-friendly spaces and plan for busy or loud events
    \item Support interests and strengths; they are assets, not problems
    \item Look after sleep, diet, activity and general health
    \item Take time to understand the person's communication, and teach alternatives if needed
    \item Get support for carers: counselling, respite and parent groups
    \item Be sceptical of miracle "cures"; ask a trusted clinician before trying anything unusual
\end{itemize}

\section*{Sources and Further Reading}
The following reputable organisations provide reliable, current information on treatments and supports for autism.

\begin{itemize}
    \item National Health Service. \textit{Information on autism, including therapies and support.} https://www.nhs.uk/conditions/autism/
    \item National Institute for Health and Care Excellence. \textit{Clinical guidance on the management and support of people with autism.} https://www.nice.org.uk/guidance/ng128
    \item Autism Spectrum Australia (Aspect). \textit{Services, programs and information for autistic people and families.} https://www.aspect.org.au
    \item Raising Children Network (Australia). \textit{Information on therapies and support for autistic children.} https://raisingchildren.net.au
    \item Centers for Disease Control and Prevention. \textit{Information on autism treatment and intervention.} https://www.cdc.gov/autism/
    \item National Disability Insurance Scheme. \textit{Information about funded supports for people with autism.} https://www.ndis.gov.au
\end{itemize}
